% from ega2-1-fr.tex
\oldpage[II]{21}
Soient $S$, $S'$ deux anneaux gradués~; un homomorphisme d'\emph{anneaux
gradués} $\varphi:S\to S'$ est un homomorphisme d'anneaux tel que
$\varphi(S_n)\subset S'_n$, pour tout $n\in\mathbf{Z}$ (autrement dit,
$\varphi$ doit être un homomorphisme \emph{de degré $0$} de
$\mathbf{Z}$-modules gradués). La donnée d'un tel homomorphisme définit sur
$S'$ une structure de $S$-module \emph{gradué}~; muni de cette structure et
de sa structure d'anneau gradué, on dit que $S'$ est une \emph{$S$-algèbre
graduée}.

Si alors $M$ est un $S$-module gradué, le produit tensoriel $M\otimes_S S'$
de $S$-modules gradués est muni de façon naturelle d'une structure de
$S'$-module \emph{gradué}, la graduation étant celle définie plus haut.
\end{env}

\begin{lemma}[2.1.3]
\label{II.2.1.3-fr}
Soit $S$ un anneau gradué à degrés positifs. Pour qu'une partie $E$ de $S_+$
formée d'éléments homogènes engendre $S_+$ en tant que $S$-module, il faut et
il suffit que $E$ engendre $S$ en tant que $S_0$-algèbre.
\end{lemma}

\begin{proof}
La condition est évidemment suffisante~; montrons qu'elle est nécessaire.
Soit $E_n$ (resp. $E^n$) l'ensemble des éléments de $E$ de degré égal à $n$
(resp. $\leq n$)~; il suffira de montrer, par récurrence sur $n>0$, que $S_n$
est le $S_0$-module engendré par les éléments de degré $n$ qui sont des
produits d'éléments de $E^n$. Or, cela est évident pour $n=1$ en vertu de
l'hypothèse~; cette dernière montre en outre que
$S_n=\sum_{p=0}^{n-1}S_pE_{n-p}$, et le raisonnement par récurrence est alors
immédiat.
\end{proof}

\begin{corollary}[2.1.4]
\label{II.2.1.4-fr}
Pour que $S_+$ soit un idéal de type fini, il faut et il suffit que $S$ soit
une $S_0$-algèbre de type fini.
\end{corollary}

\begin{proof}
On peut en effet toujours supposer qu'un système fini de générateurs de la
$S_0$-algèbre $S$ (resp. du $S$-idéal $S_+$) est formé d'éléments homogènes,
en remplaçant chacun des générateurs considérés par ses composants homogènes.
\end{proof}

\begin{corollary}[2.1.5]
\label{II.2.1.5-fr}
Pour que $S$ soit noethérien, il faut et il suffit que $S_0$ soit noethérien
et que $S$ soit une $S_0$-algèbre de type fini.
\end{corollary}

\begin{proof}
La condition est évidemment suffisante~; elle est nécessaire, puisque $S_0$
est isomorphe à $S/S_+$ et que $S_+$ doit être un idéal de type fini
(\hyperref[II.2.1.4-fr]{2.1.4}).
\end{proof}

\begin{lemma}[2.1.6]
\label{II.2.1.6-fr}
Soit $S$ un anneau gradué à degrés positifs, qui soit une $S_0$-algèbre de
type fini. Soit $M$ un $S$-module gradué de type fini. Alors~:
\begin{enumerate}
  \item[\rm(i)] Les $M_n$ sont des $S_0$-modules de type fini, et il existe
    un entier $n_0$ tel que $M_n=0$ pour $n\leq n_0$.
  \item[\rm(ii)] Il existe un entier $n_1$ et un entier $h>0$ tel que, pour
    tout entier $n\geq n_1$, on ait $M_{n+h}=S_hM_n$.
  \item[\rm(iii)] Pour tout couple d'entiers $(d,k)$ tels que $d>0$,
    $0\leq k\leq d-1$, $M^{(d,k)}$ est un $S^{(d)}$-module de type fini.
  \item[\rm(iv)] Pour tout entier $d>0$, $S^{(d)}$ est une $S_0$-algèbre de
    type fini.
  \item[\rm(v)] Il existe un entier $h>0$ tel que $S_{mh}=(S_h)^m$ pour tout
    $m>0$.
  \item[\rm(vi)] Pour tout entier $n>0$, il existe un entier $m_0$ tel que
    $S_m\subset S_+^n$ pour tout $m\geq m_0$.
\end{enumerate}
\end{lemma}

\begin{proof}
On peut supposer $S$ engendré (en tant que $S_0$-algèbre) par des éléments
homogènes $f_i$, de degré $h_i$ ($1\leq i\leq r$), et $M$ engendré (en tant
que $S$-module) par des éléments homogènes $x_j$ de degré $k_j$
($1\leq j\leq s$). Il est clair que $M_n$ est formé des combinaisons
